% Options for packages loaded elsewhere
\PassOptionsToPackage{unicode}{hyperref}
\PassOptionsToPackage{hyphens}{url}
\PassOptionsToPackage{dvipsnames,svgnames,x11names}{xcolor}
\documentclass[
  10pt,
  fontset=fandol]{ctexart}
\usepackage{xcolor}
\usepackage[margin=16mm]{geometry}
\usepackage{amsmath,amssymb}
\setcounter{secnumdepth}{-\maxdimen} % remove section numbering
\usepackage{iftex}
\ifPDFTeX
  \usepackage[T1]{fontenc}
  \usepackage[utf8]{inputenc}
  \usepackage{textcomp} % provide euro and other symbols
\else % if luatex or xetex
  \usepackage{unicode-math} % this also loads fontspec
  \defaultfontfeatures{Scale=MatchLowercase}
  \defaultfontfeatures[\rmfamily]{Ligatures=TeX,Scale=1}
\fi
\usepackage{lmodern}
\ifPDFTeX\else
  % xetex/luatex font selection
\fi
% Use upquote if available, for straight quotes in verbatim environments
\IfFileExists{upquote.sty}{\usepackage{upquote}}{}
\IfFileExists{microtype.sty}{% use microtype if available
  \usepackage[]{microtype}
  \UseMicrotypeSet[protrusion]{basicmath} % disable protrusion for tt fonts
}{}
\makeatletter
\@ifundefined{KOMAClassName}{% if non-KOMA class
  \IfFileExists{parskip.sty}{%
    \usepackage{parskip}
  }{% else
    \setlength{\parindent}{0pt}
    \setlength{\parskip}{6pt plus 2pt minus 1pt}}
}{% if KOMA class
  \KOMAoptions{parskip=half}}
\makeatother
\setlength{\emergencystretch}{3em} % prevent overfull lines
\providecommand{\tightlist}{%
  \setlength{\itemsep}{0pt}\setlength{\parskip}{0pt}}
\usepackage{amsmath,amssymb,mathtools,needspace}
\xeCJKDeclareCharClass{CJK}{"2032,"0400->"04FF,"2160->"216B,"2460->"2473,"25A0,"25C7,"25CB,"25CF}
\IfFontExistsTF{Arial Unicode MS}{\xeCJKsetup{AutoFallBack=true}\setCJKfallbackfamilyfont{\CJKrmdefault}{Arial Unicode MS}}{}
\setcounter{secnumdepth}{0}
\setlength{\emergencystretch}{2em}
\allowdisplaybreaks
\usepackage{bookmark}
\IfFileExists{xurl.sty}{\usepackage{xurl}}{} % add URL line breaks if available
\urlstyle{same}
\hypersetup{
  pdftitle={分段三次 Hermite 插值},
  colorlinks=true,
  linkcolor={Maroon},
  filecolor={Maroon},
  citecolor={Blue},
  urlcolor={Blue},
  pdfcreator={LaTeX via pandoc}}

\title{分段三次 Hermite 插值}
\author{}
\date{}

\begin{document}
\maketitle

\subsubsection{2.7.3 分段三次 Hermite 插值}\label{ux5206ux6bb5ux4e09ux6b21-hermite-ux63d2ux503c}

分段线性插值函数 \(I_h(x)\) 的导数是间断的，若在节点 \(x_k\)（\(k=0,1,\cdots,n\)）上除已知函数值 \(f_k\) 外还给出导数值 \(f'_k=m_k\)（\(k=0,1,\cdots,n\)），就可构造一个导数连续的分段插值函数 \(I_h(x)\)，它满足如下条件：

（1）\(I_h(x)\in C^1[a,b]\)（\(C^1[a,b]\) 代表区间 \([a,b]\) 上一阶导数连续的函数集合）；

（2）\(I_h(x_k)=f_k,I'_h(x_k)=f'_k\)（\(k=0,1,\cdots,n\)）；

（3）\(I_h(x)\) 在每个小区间 \([x_k,x_{k+1}]\) 上是三次多项式。

根据两点三次插值多项式（2.6.10）可知，\(I_h(x)\) 在区间 \([x_k,x_{k+1}]\) 上的表达式为

\[\begin{aligned}
I_h(x)={}&\left(\frac{x-x_{k+1}}{x_k-x_{k+1}}\right)^2\left(1+2\frac{x-x_k}{x_{k+1}-x_k}\right)f_k+\left(\frac{x-x_k}{x_{k+1}-x_k}\right)^2\left(1+2\frac{x-x_{k+1}}{x_k-x_{k+1}}\right)f_{k+1}\\
&+\left(\frac{x-x_{k+1}}{x_k-x_{k+1}}\right)^2(x-x_k)f'_k+\left(\frac{x-x_k}{x_{k+1}-x_k}\right)^2(x-x_{k+1})f'_{k+1}.
\end{aligned}\tag{2.7.5}\]

若在整个区间 \([a,b]\) 上定义一组分段三次插值基函数 \(\alpha_j(x)\) 及 \(\beta_j(x)\)（\(j=0,\)

\begin{quote}
续页：上句的下标范围及公式在下一页续完。
\end{quote}

\href{../../source-images/pdf-048.jpeg}{查看原页}

\begin{quote}
本页接上页 2.7.3 正文；开头续完下标范围。
\end{quote}

\(1,\cdots,n\)），则 \(I_h(x)\) 可表示为

\[I_h(x)=\sum_{j=0}^n[f_j\alpha_j(x)+f'_j\beta_j(x)],\tag{2.7.6}\]

其中 \(\alpha_j(x),\beta_j(x)\) 依式（2.6.8）、式（2.6.9）分别表示为

\[\alpha_j(x)=\begin{cases}
\displaystyle\left(\frac{x-x_{j-1}}{x_j-x_{j-1}}\right)^2\left(1+2\frac{x-x_j}{x_{j-1}-x_j}\right),&x_{j-1}\leqslant x\leqslant x_j\quad(j=0\text{ 略去}),\\[6pt]
\displaystyle\left(\frac{x-x_{j+1}}{x_j-x_{j+1}}\right)^2\left(1+2\frac{x-x_j}{x_{j+1}-x_j}\right),&x_j\leqslant x\leqslant x_{j+1}\quad(j=n\text{ 略去}),\\[6pt]
0,&\text{其他},
\end{cases}\tag{2.7.7}\]

\[\beta_j(x)=\begin{cases}
\displaystyle\left(\frac{x-x_{j-1}}{x_j-x_{j-1}}\right)^2(x-x_j),&x_{j-1}\leqslant x\leqslant x_j\quad(j=0\text{ 略去}),\\[6pt]
\displaystyle\left(\frac{x-x_{j+1}}{x_j-x_{j+1}}\right)^2(x-x_j),&x_j\leqslant x\leqslant x_{j+1}\quad(j=n\text{ 略去}),\\[6pt]
0,&\text{其他}.
\end{cases}\tag{2.7.8}\]

由于 \(\alpha_j(x),\beta_j(x)\) 的局部非零性质，当 \(x\in[x_k,x_{k+1}]\) 时，只有 \(\alpha_k(x),\alpha_{k+1}(x),\beta_k(x),\beta_{k+1}(x)\) 不为零，于是式（2.7.6）的 \(I_h(x)\) 可表示为

\[I_h(x)=f_k\alpha_k(x)+f_{k+1}\alpha_{k+1}(x)+f'_k\beta_k(x)+f'_{k+1}\beta_{k+1}(x)\quad(x_k\leqslant x\leqslant x_{k+1}).\tag{2.7.9}\]

为了研究 \(I_h(x)\) 的收敛性，由式（2.7.7）及式（2.7.8）直接得估计式

\[0\leqslant\alpha_j(x)\leqslant1,\tag{2.7.10}\]

\[\begin{cases}
\displaystyle|\beta_k(x)|\leqslant\frac{4}{27}h_k,\\[6pt]
\displaystyle|\beta_{k+1}(x)|\leqslant\frac{4}{27}h_k.
\end{cases}\tag{2.7.11}\]

此外，当 \(f(x)\) 是分段三次多项式时，\(f(x)\) 的插值多项式 \(I_h(x)\) 就是它本身。例如，当 \(f(x)=1\) 时就有 \(\sum_{j=0}^n\alpha_j(x)=1\)。当 \(x\in[x_k,x_{k+1}]\) 时，得

\[\alpha_k(x)+\alpha_{k+1}(x)=1.\tag{2.7.12}\]

由式（2.7.9）～（2.7.12），当 \(x\in[x_k,x_{k+1}]\) 时还可得

\[\begin{aligned}
|f(x)-I_h(x)|&\leqslant\alpha_k(x)|f(x)-f_k|+\alpha_{k+1}(x)|f(x)-f_{k+1}|\\
&\quad+\frac{4}{27}h_k[|f'_k|+|f'_{k+1}|]\\
&\leqslant[\alpha_k(x)+\alpha_{k+1}(x)]\omega(h)+\frac{8}{27}h\max\{|f'_k|,|f'_{k+1}|\},
\end{aligned}\]

\[|f(x)-I_h(x)|\leqslant\omega(h)+\frac{8}{27}h\max_{0\leqslant k\leqslant n}|f'_k|\tag{2.7.13}\]

对于 \(x\in[a,b]\) 成立，其中 \(\omega(h)\) 是 \(f(x)\) 在 \([a,b]\) 上的连续模。因此，当 \(f(x)\in C[a,b]\)

\begin{quote}
续页：收敛性论述在下一页续完。
\end{quote}

\href{../../source-images/pdf-049.jpeg}{查看原页}

\begin{quote}
本页接上页 2.7.3 正文。
\end{quote}

时，

\[\lim_{h\to0}I_h(x)=f(x).\]

这就证明了 \(I_h(x)\) 在 \([a,b]\) 上一致收敛到 \(f(x)\)。

\begin{center}\rule{0.5\linewidth}{0.5pt}\end{center}

\subsection{相关原页校注（agent 补充）}\label{ux76f8ux5173ux539fux9875ux6821ux6ce8agent-ux8865ux5145}

以下校注来自本节覆盖的原页，可能也涉及同页相邻小节；原文照录与校正说明分开保留。

\subsubsection{原书 PDF 页 49 的校注}\label{ux539fux4e66-pdf-ux9875-49-ux7684ux6821ux6ce8}

\href{../pages/pdf-049.md}{完整原页转录} · \href{../../source-images/pdf-049.jpeg}{原页图像}

校注记录：ch02b-E002。

\begin{quote}
\textbf{校注（agent 补充；原书论证条件疑漏）}：本页开头承接上一页``当 \(f(x)\in C[a,b]\) 时''即宣称分段三次 Hermite 插值一致收敛，原文已保留。由（2.7.13）实际还需控制 \(h\max_k|f'_k|\to0\)；仅有连续性既不保证节点导数存在，也不能控制随加密节点变化的导数值。比如加上 \(f\in C^1[a,b]\) 就足以使节点导数一致有界，从该估计推出结论。本注是对原文论证条件的补充，不属于教材正文。
\end{quote}

\subsection{本节覆盖原页的校注记录}\label{ux672cux8282ux8986ux76d6ux539fux9875ux7684ux6821ux6ce8ux8bb0ux5f55}

下列记录按原页关联，含同页相邻小节的校注；计算时同时读取上文校注和完整原页。

\begin{itemize}
\tightlist
\item
  \texttt{ch02b-E002}：原书 PDF 页 49
\end{itemize}

\end{document}
